% --- Patron: Aveh, the Rider Behind the Storm (Freedom & Erasure) ---
\subsection{Aveh, the Rider Behind the Storm (Freedom \& Erasure)}

\textit{Lore.} Aveh is the faceless rider at the horizon's edge, neither man nor woman, both and neither. The Ykrul call them \textit{Aveh}, a liminal spirit beyond hearth and vow, a storm-shadow that offers liberty at the cost of belonging. To travel with Aveh is to taste unchained freedom---and to risk being erased by the storm that follows.

\begin{quote}
``I am the track the storm devours. I am the freedom you fear and the silence you crave. To ride with me is to be unbound, and to be unbound is to vanish.''
\end{quote}

\subsubsection*{Rites of the Rider}

\paragraph{Storm-Step (Low, 4 XP)} \hfill \\
\emph{Scene; Self; Resist only.} \\
\textbf{Materials:} A breath cast into the wind. \\
\textbf{Effect:} Slip free of notice or restraint. Gain \textbf{+1 die} on one escape or evasion roll this scene. \\
\textbf{Push It:} Instead gain \textbf{+2 dice}, but mark \textbf{1 SB (Spades)} as your absence unsettles the scene. \\
\emph{Requires: Familiar (\textit{Invoke:} 1 Boon).}

\paragraph{Exile's Banner (Low, 5 XP)} \hfill \\
\emph{Scene; Near; No.} \\
\textbf{Materials:} A line drawn in dust, ash, or stormwater. \\
\textbf{Effect:} Mark a target as "outside." The target suffers \textbf{-1 die} on all social rolls for the scene. \\
\textbf{Push It:} The target is fully estranged; allies treat them as forgotten, but you immediately advance your Corruption clock 1 segment. \\
\emph{Requires: Familiar (\textit{Invoke:} 1 Boon).}

\paragraph{Erase the Road (Standard, 8 XP)} \hfill \\
\emph{Scene; Location; No.} \\
\textbf{Materials:} A threshold or roadway under open sky. \\
\textbf{Effect:} All attempts to follow you this session suffer \textbf{-2 dice}. In exchange, erase part of your own path: lose a detail (a memory, a name) chosen by the GM. \\
\textbf{Push It:} Expand the effect to an entire group, but mark \textbf{1 SB (Clubs)} as reality warps around the erasure. \\
\emph{Requires: Familiar + Codex (\textit{Invoke:} 1 Boon).}

\paragraph{The Rider's Mark (Standard, 9 XP)} \hfill \\
\emph{Scene; Touch; No.} \\
\textbf{Materials:} Stormwater or dust rubbed into skin. \\
\textbf{Effect:} A companion gains \textbf{+2 dice} to resist capture or coercion, but others forget them more easily. \\
\textbf{Push It:} Grant \textbf{+3 dice} instead, but mark \textbf{1 SB (Diamonds)} as the Rider's storm distorts bonds. \\
\emph{Requires: Familiar + Codex (\textit{Invoke:} 1 Boon).}

\paragraph{Storm's Refuge (High, 12 XP)} \hfill \\
\emph{Scene; Zone; No.} \\
\textbf{Materials:} A banner or cloak raised in storm. \\
\textbf{Effect:} Your group is unseen or forgotten by pursuers for the scene. At the end, all present must mark 1 Corruption or forget one bond. \\
\textbf{Push It:} Extend protection for an entire session, but mark \textbf{2 SB (Hearts)}. \\
\emph{Requires: Familiar + Codex + Tier III (\textit{Invoke:} 2 Boons).} \\
\emph{Obligation:} 6 segments.

\paragraph{The Horizon Devours (High, 14 XP)} \hfill \\
\emph{Extended; Zone; No.} \\
\textbf{Materials:} A circle of ash or dust drawn at a boundary. \\
\textbf{Effect:} All thresholds in a zone collapse. In exchange, something of you is erased---your reflection, your name, or your presence in memory. \\
\textbf{Push It:} The storm devours even time; mark \textbf{3 SB (Spades)} as history forgets you. \\
\emph{Requires: Familiar + Codex + Tier III (\textit{Invoke:} 2 Boons).} \\
\emph{Obligation:} 8 segments.

\subsubsection*{Cantors \& Cults of Aveh}

\textbf{The Cantor's Song.} A Cantor of Aveh does not sing to be heard, but to be forgotten. Their art is the \textit{empty note}, the melody that leaves a silence where a memory once lived. They are sought by those who need to disappear---spies, exiles, the grief-stricken---and feared by those who want to be remembered.

\textbf{The Cult of the Unbannered.} Aveh has no temples, only crossroads. Her cults are small, mobile, and secretive. A ritual of the Unbannered might involve each member writing their name on a slip of paper, then passing the slips in a circle until no one remembers whose name they hold. They are less a congregation and more a mutual aid society for the erased, offering shelter, forged papers, and a shared promise: \textit{``We were never here.''}

\subsubsection*{Witchcraft of Aveh (Threshold)} \hfill \\
A witch who walks with Aveh is a \textit{path-cutter}, a specialist in hidden roads and forgotten boundaries. Their power lies in \textit{omission}: removing a door, erasing a trail, or unlatching the memory of a promise.

\textbf{The Witch's Tool: The Wind-Key.} A small iron whistle, cold to the touch even in summer. When blown at a locked threshold---a door, a city gate, a sealed letter---the wind answers, and the lock opens as if it had never been closed. Once per session.

\textbf{A Hedge Gift: Unlatched (2 XP).} Once per scene, you may declare a mundane lock, latch, or simple binding to be ``unlatched'' and open freely.

\subsubsection*{Aveh's Corruption Table}

% Requires \usepackage{tabularray} and \usepackage{xcolor} in preamble
\begin{tblr}{
  colspec = {Q[c,1cm] X[2.5] X[2.5]},
  rowhead = 1,
  row{odd} = {bg=gray!10},
  row{1} = {bg=gray!30, font=\bfseries},
  hlines,
}
Tier & Benefit & Cost \\
1 & \textbf{Storm-sense:} +1 die to Notice rolls outdoors. & \textbf{Voice falters:} Others forget your words more easily. \\
2 & \textbf{Fleetness:} +1 die on escape or chase rolls. & \textbf{Shadow lags:} Mark 1 Fatigue when first seen in a scene. \\
3 & \textbf{Erasure's Mercy:} Once/session, negate 1 Harm or capture. & \textbf{Memory slips:} Allies forget one fact or bond tied to you. \\
4 & \textbf{Rider's Endurance:} Resist 1 Condition (Fear, Fatigue, Bind) once/session. & \textbf{Reflections falter:} You do not appear clearly in mirrors. \\
5 & \textbf{Storm-banner:} +1 die to inspire rebellion or defiance. & \textbf{Exile's weight:} Each Push of a Rite marks 1 permanent Corruption. \\
6+ & \textbf{Horizon's Claim:} You may ride unseen across any threshold. & \textbf{Self-erasure:} On social rolls, roll twice and take the worst result. \\
\end{tblr}